% --- 自動計數器設定 ---
\newtotcounter{qcount} \regtotcounter{qcount}
\newtotcounter{sccount} \regtotcounter{sccount} 
\newtotcounter{mccount} \regtotcounter{mccount} 
\newtotcounter{wtscore} \regtotcounter{wtscore} % 非選題總分計數器
\newcommand{\scscore}{3}
\newcommand{\mcscore}{4}

% --- 抬頭設定 ---
\newcommand{\examheader}{
    \begin{center}
        \section*{第一章節複習卷}
    \end{center}
}

% --- 基礎縮排設定 ---
\newcommand{\choiceindent}{0.85cm}

% --- 1. 防分頁容器 (qblock) ---
\newif\ifqblockpassthrough

\newenvironment{qblock}{%
    \ifqblockpassthrough
    \else
        \par\vspace{0.5em}
        \noindent
        \begin{minipage}{\linewidth}
    \fi
}{%
    \ifqblockpassthrough
    \else
        \end{minipage}
        \par\vspace{1em}
    \fi
}
% --- 2. 單選題指令 ---
\newcommand{\schoices}[7]{% 
    \begin{qblock}
        \stepcounter{qcount}\stepcounter{sccount}%
        \begin{enumerate}[label=\textbf{\theqcount.}, leftmargin=\choiceindent, labelsep=0.2cm, nosep]
            \item #1
        \end{enumerate}\vspace{0.2em}
        \ifstrequal{#2}{h}{
            % 水平模式：利用 \leftskip 讓整段文字的左邊界鎖定在 \choiceindent
            {\leftskip=\choiceindent \noindent
            (A) #3 \hspace{0.5cm} (B) #4 \hspace{0.5cm} (C) #5 \hspace{0.5cm} (D) #6 \hspace{0.5cm} (E) #7 \par}
        }{
            \ifstrequal{#2}{m}{
                \noindent\hspace*{\choiceindent}%
                \begin{tabular}[t]{@{}l @{\hspace{0.5cm}} l @{\hspace{0.5cm}} l @{\hspace{0.5cm}} l @{\hspace{0.5cm}} l}
                    (A) #3 & (B) #4 & (C) #5 & (D) #6 & (E) #7
                \end{tabular}\par
            }{
                \begin{enumerate}[(A), topsep=0.3em, itemsep=0.5em, parsep=0pt, leftmargin=*, labelindent=\choiceindent, labelsep=0.25cm, align=left] 
                    \item #3 \item #4 \item #5 \item #6 \item #7
                \end{enumerate}
            }
        }
    \end{qblock}
}

% --- 3. 複選題指令 ---
\newcommand{\mchoices}[8]{% 
    \begin{qblock}
        \stepcounter{qcount}\stepcounter{mccount}%
        \begin{enumerate}[label=\textbf{\theqcount.}, leftmargin=\choiceindent, labelsep=0.2cm, nosep]
            \item #1 \textbf{(應選 #2 項)}
        \end{enumerate}\vspace{0.2em}
        \ifstrequal{#3}{h}{
            % 水平模式同步修改
            {\leftskip=\choiceindent \noindent
            (A) #4 \hspace{0.5cm} (B) #5 \hspace{0.5cm} (C) #6 \hspace{0.5cm} (D) #7 \hspace{0.5cm} (E) #8 \par}
        }{
            \ifstrequal{#3}{m}{
                \noindent\hspace*{\choiceindent}%
                \begin{tabular}[t]{@{}l @{\hspace{0.5cm}} l @{\hspace{0.5cm}} l @{\hspace{0.5cm}} l @{\hspace{0.5cm}} l}
                    (A) #4 & (B) #5 & (C) #6 & (D) #7 & (E) #8
                \end{tabular}\par
            }{
                \begin{enumerate}[(A), topsep=0.3em, itemsep=0.5em, parsep=0pt, leftmargin=*, labelindent=\choiceindent, labelsep=0.25cm, align=left] 
                    \item #4 \item #5 \item #6 \item #7 \item #8
                \end{enumerate}
            }
        }
    \end{qblock}
}

% --- 側邊圖片包裝指令 ---
% #1: 題目內容
% #2: 圖片檔名
% #3: 圖片佔頁面寬度比例
% #4: 模式選擇: w (文繞圖) / n (原版左右並排)
\newlength{\sideimgwidth}

\newcommand{\sideimg}[4]{%
    \begin{qblock}
        \ifstrequal{#4}{w}{%
            \setlength{\sideimgwidth}{#3\linewidth}%

            \begin{wrapstuff}[
                r,
                width=\sideimgwidth,
                hsep=0.8em,
                voffset=-0.5em
            ]
                \includegraphics[
                    width=\sideimgwidth
                ]{#2}
            \end{wrapstuff}

            \qblockpassthroughtrue
            #1
            \qblockpassthroughfalse

            \wrapstuffclear
        }{%
            \noindent
            \begin{minipage}[t]{%
                \dimexpr0.91\linewidth-#3\linewidth\relax
            }
                \vspace{-0.5em}
                \qblockpassthroughtrue
                #1
                \qblockpassthroughfalse
            \end{minipage}
            \hfill
            \begin{minipage}[t]{#3\linewidth}
                \includegraphics[
                    width=\linewidth,
                    valign=t
                ]{#2}
            \end{minipage}
        }
    \end{qblock}
}

% --- 4. 非選題指令 ---
% (1) 單一非選題 (無子題): \writing{該題總分}{題幹敘述}
\newcommand{\writing}[2]{%
    \begin{qblock}
        \stepcounter{qcount}% 
        \addtocounter{wtscore}{#1}% 將此題分數累加至「非選總分」
        \begin{enumerate}[label=\textbf{\theqcount.}, leftmargin=\choiceindent, labelsep=0.2cm, nosep]
            \item #2 \textbf{(共 #1 分)}
        \end{enumerate}\vspace{0.2em}
    \end{qblock}
}

% (2) 多子題非選題: \writings{該題總分}{題幹敘述}{子題內容}
\newcommand{\writings}[3]{%
    \begin{qblock}
        \stepcounter{qcount}% 
        \addtocounter{wtscore}{#1}% 
        \begin{enumerate}[label=\textbf{\theqcount.}, leftmargin=\choiceindent, labelsep=0.2cm, nosep]
            \item #2 \textbf{(共 #1 分)}
        \end{enumerate}\vspace{0.2em}
        \ifblank{#3}{}{
            \begin{enumerate}[label=(\arabic*), leftmargin=\dimexpr\choiceindent+0.8cm\relax, labelsep=0.2cm, topsep=0.2em, itemsep=0.5em]
                #3
            \end{enumerate}
        }
    \end{qblock}
}

% (3) 子題專用指令: \subq{配分}{敘述}
\newcommand{\subq}[2]{%
    \item #2 \textbf{(#1 分)}
}

% --- 5. 選項排版 ---
\newcommand{\hchoices}[5]{% 水平排列
    \noindent\hspace*{\choiceindent}% 
    (A) #1 \hfill (B) #2 \hfill (C) #3 \hfill (D) #4 \hfill (E) #5 \par
}

\newcommand{\vchoices}[5]{% 垂直排列
    \begin{enumerate}[(A), nosep, leftmargin=*, labelindent=\choiceindent, labelsep=0.25cm, align=left] 
        \item #1 \item #2 \item #3 \item #4 \item #5
    \end{enumerate}
}

% ---------- 自動答案表格 ----------

% 每格寬度
\newlength{\answercellwidth}
\setlength{\answercellwidth}{1.5cm}

% 每格高度倍率
\newcommand{\answerarraystretch}{1.8}

\ExplSyntaxOn

\clist_new:N \l__autoanswer_clist
\int_new:N \l__autoanswer_columns_int

% 產生題號列
\cs_new_protected:Npn \__autoanswer_number_row:n #1
{
    \int_step_inline:nnnn
        {0}
        {1}
        {\l__autoanswer_columns_int - 1}
    {
        \int_compare:nT
            {
                #1 + ##1
                <=
                \clist_count:N \l__autoanswer_clist
            }
            {
                \int_eval:n {#1 + ##1}
            }

        \int_compare:nNnT
            {##1}
            <
            {\l__autoanswer_columns_int - 1}
            {&}
    }

    \tabularnewline
    \hline
}

% 產生答案列
\cs_new_protected:Npn \__autoanswer_answer_row:n #1
{
    \int_step_inline:nnnn
        {0}
        {1}
        {\l__autoanswer_columns_int - 1}
    {
        \int_compare:nT
            {
                #1 + ##1
                <=
                \clist_count:N \l__autoanswer_clist
            }
            {
                \clist_item:Nn
                    \l__autoanswer_clist
                    {#1 + ##1}
            }

        \int_compare:nNnT
            {##1}
            <
            {\l__autoanswer_columns_int - 1}
            {&}
    }
}

% 建立完整表格
\cs_new_protected:Npn \__autoanswer_table:nn #1#2
{
    \group_begin:

    % #1：每列題數
    % #2：答案清單
    \int_set:Nn \l__autoanswer_columns_int {#1}
    \clist_set:Nn \l__autoanswer_clist {#2}

    \setlength{\tabcolsep}{0pt}
    \renewcommand{\arraystretch}{\answerarraystretch}

    \begin{tabular}{
        |*{\int_use:N \l__autoanswer_columns_int}{
            >{\centering\arraybackslash}
            m{\answercellwidth}|
        }
    }
        \hline

        \int_step_inline:nnnn
            {1}
            {\l__autoanswer_columns_int}
            {\clist_count:N \l__autoanswer_clist}
        {
            % 從第二組開始，先關閉上一個答案列
            \int_compare:nNnT
                {##1}
                >
                {1}
                {
                    \tabularnewline
                    \hline
                }

            \__autoanswer_number_row:n {##1}
            \__autoanswer_answer_row:n {##1}
        }

        % 關閉最後一列
        \tabularnewline
        \hline
    \end{tabular}

    \group_end:
}

% 預設每列 5 題
\NewDocumentCommand{\autoanswer}{O{5}m}
{
    \__autoanswer_table:nn {#1} {#2}
}

\ExplSyntaxOff
